\appendix
\chapter{Artifacts and Reproduction}
\label{app:artifacts}
\sloppy

\section{Artifact Location}

The development repository is available at
\begin{center}
\url{https://github.com/anton-fadeenkov/stella-lsp-extensions}.
\end{center}

The submitted thesis archive contains the exact source revision used for the screenshots and manual scenarios. Because the development repository may change, evaluation should use the submitted archive or a release tag associated with the thesis.

\section{Build and Launch}

The following commands are executed from the repository root:

\begin{enumerate}
    \item Install dependencies with \texttt{npm install}.
    \item Build the extension with \texttt{npm run build}.
    \item Open the repository in Visual Studio Code.
    \item Start an Extension Development Host and open one of the validation \texttt{.stella} files.
\end{enumerate}

The commands above reproduce the implementation used for the manual scenarios. The Visual Studio Code commands exposed by the artifact are \texttt{Stella: Show Type Derivation}, \texttt{Stella: Show Type Error Trace}, \texttt{Stella: Show Type Dependency Graph}, and \texttt{Stella: Show AST (JSON)}.

\section{Implementation Map}

The main implementation artifacts are listed below.

\begin{itemize}
    \item \textbf{AST visualization:}
    client: \path{src/extension/syntax-tree.ts};
    server: \path{src/language/lsp/syntax-request.ts}.

    \item \textbf{Scope and context:}
    client: \path{src/extension/scope-view.ts};
    server: \path{src/language/lsp/scope-request.ts}.

    \item \textbf{Type derivation:}
    client: \path{src/extension/type-derivation-panel.ts};
    server: \path{src/language/lsp/type-derivation-request.ts}.

    \item \textbf{Type-error trace:}
    client: \path{src/extension/type-error-trace-panel.ts};
    server: \path{src/language/lsp/type-error-trace-request.ts}.

    \item \textbf{Type dependency graph:}
    client: \path{src/extension/type-dependency-graph-panel.ts};
    server: \path{src/language/lsp/type-dependency-graph-request.ts}.

    \item \textbf{Shared response structures:}
    \path{src/shared/lsp} and \path{src/shared/analysis}.
\end{itemize}

\section{Validation Inputs}

The validation files are included with the artifact:

\begin{itemize}
    \item dependency-graph programs: \path{examples/dependency-graph};
    \item expected dependency nodes and edges: \path{examples/type-dependency-graph-samples.md};
    \item error-trace inputs and expected rules: \path{examples/error-trace-samples.md}.
\end{itemize}

To repeat a manual scenario:

\begin{enumerate}
    \item Open a listed \texttt{.stella} file in the Extension Development Host.
    \item Place the cursor on the expression named in the scenario description.
    \item Run the corresponding \texttt{Stella: Show ...} command.
    \item Compare nodes, type labels, edge labels, and conflict markers with the documented expectation.
    \item Click visual elements and verify that the expected source ranges are revealed.
\end{enumerate}

The screenshots used in Chapter~\ref{chap:implementation} are stored in
\path{Anton_Fadeenkov/figs}.

